%!TEX root = ../QuantumNoiseAndDecoherence.tex
% ──────────────────────────────────────────────────────────────
%  Lecture 11 — Continuous Measurement and the Lindblad Equation
% ──────────────────────────────────────────────────────────────

% Continues Section 7 (Classical noise and stochastic calculus);
% opens Section 8 (Quantum stochastic calculus).

\subsection{From repeated measurements to continuous evolution}%
\label{sec:repeated-to-continuous}

We now carry out the programme announced in the previous lecture:
starting from $n$~discrete von~Neumann measurements of position, we
take $n \to \infty$ and $\tau \to 0$ simultaneously, obtaining a
Lindblad master equation that models \emph{continuous measurement}.

\begin{intuition}[Decoherence as information flow]
  A system loses coherence when the environment gains information
  about it.  Information gain builds up entanglement between system
  and environment.  Continuous measurement provides a transparent
  model for this process.
\end{intuition}

\subsubsection{Hilbert space and Hamiltonian}

Recall the setup: the system is a continuous quantum variable with
canonical operators $\hat{x}_S$, $\hat{p}_S$ satisfying
$[\hat{x}_S, \hat{p}_S] = i$, and we introduce $n$~independent
meter systems with canonical pairs $(\hat{x}_j, \hat{p}_j)$.
The total Hilbert space is
\begin{equation}\label{eq:cm-hilb}
  \hilb_{\mathrm{tot}}
  = \hilb_{M_1} \tp \cdots \tp \hilb_{M_n} \tp \hilb_S\,,
  \qquad \hilb_{M_j} = L^2(\R)\,.
\end{equation}
The time-dependent Hamiltonian modelling $n$~instantaneous
measurements at times $t_j = j\tau$ is
\begin{equation}\label{eq:cm-hamiltonian}
  \eqbox{H(t) = H_0
  + \sum_{j=1}^{n} \delta(t - j\tau)\;\hat{x}_S \tp \hat{p}_j}\,,
\end{equation}
where $H_0$ is the free system Hamiltonian.
Each meter is prepared in a Gaussian state of width~$\sqrt{\sigma}$:
\begin{equation}\label{eq:meter-gaussian}
  g_j(x_j) = \braket{x_j}{g_j}
  = \frac{1}{(\pi\sigma)^{1/4}}\,
    \exp\!\Bigl(-\frac{x_j^2}{2\sigma}\Bigr)\,,
\end{equation}
so the total initial state is
$\rho(0) = \ketbra{g_1}{g_1} \tp \cdots \tp \ketbra{g_n}{g_n}
\tp \rho_S$.

\begin{remark}
  The parameter $\sigma$ controls measurement resolution: a narrow
  Gaussian ($\sigma \to 0$) approximates a projective measurement;
  a wide one ($\sigma \to \infty$) yields almost no information.
\end{remark}

% ──────────────────────────────────────────────────────────────
\subsection{Selective and non-selective state update}%
\label{sec:state-update}

Write $\rho_S(t_j^-)$ and $\rho_S(t_j^+)$ for the system state
just before and just after the $j$-th interaction.  Integrating the
Schr\"odinger equation across the delta-function impulse gives
\begin{equation}\label{eq:post-impulse}
  \rho_S(t_j^+)
  = \tr_{M_j}\!\Bigl[
      U_j\,\bigl(\ketbra{g_j}{g_j} \tp \rho_S(t_j^-)\bigr)\,
      U_j^\adj \Bigr]\,,
  \quad U_j = e^{-i\,\hat{x}_S \tp \hat{p}_j}\,.
\end{equation}

\paragraph{Selective measurement.}
If we measure the pointer and obtain outcome $\bar{x}_j$, the
system is updated to the sub-normalised conditional state
$\tilde{\rho}_S(t_j^+ \mid \bar{x}_j)
= \hat{\Upsilon}(\bar{x}_j)\,\rho_S(t_j^-)\,
  \hat{\Upsilon}(\bar{x}_j)^\adj$,
where the \textbf{measurement operator} is
\begin{equation}\label{eq:upsilon-op}
  \eqbox{\hat{\Upsilon}(\bar{x}_j)
  = \frac{1}{(\pi\sigma)^{1/4}}\,
    \exp\!\Bigl(-\frac{(\hat{x}_S - \bar{x}_j)^2}{2\sigma}
    \Bigr)}\,.
\end{equation}
The normalised posterior is
$\rho_S(t_j^+ \mid \bar{x}_j)
= \tilde{\rho}_S / p(\bar{x}_j)$ with
$p(\bar{x}_j)
= \tr[\hat{\Upsilon}(\bar{x}_j)^2\,\rho_S(t_j^-)]$.

\paragraph{Non-selective measurement.}
If the meter outcome is discarded (or controlled by the
environment), we average over all outcomes:
\begin{equation}\label{eq:nonselective-update}
  \eqbox{\rho_S(t_j^+)
  = \int_{-\infty}^{\infty}\dd\bar{x}_j\;
    \hat{\Upsilon}(\bar{x}_j)\;\rho_S(t_j^-)\;
    \hat{\Upsilon}(\bar{x}_j)^\adj}\,.
\end{equation}
This defines a completely positive, trace-preserving map on~$\rho_S$.

\begin{intuition}[Measurement back-action]
  In quantum mechanics, information gain always entails disturbance.
  Even when the meter outcome is discarded, the interaction has
  irreversibly altered the system state---this is the physical
  origin of decoherence.
\end{intuition}

% ──────────────────────────────────────────────────────────────
\subsection{Full evolution and measurement trajectories}%
\label{sec:full-evolution}

After $j$~measurements with outcomes
$\bar{x}_1, \ldots, \bar{x}_j$, the conditional system state is
\begin{equation}\label{eq:selective-full}
  \rho_S(t_j^+ \mid \bar{x}_1, \ldots, \bar{x}_j)
  = \frac{\hat{M}_j \;\rho_S(0)\; \hat{M}_j^\adj}
         {p(\bar{x}_1, \ldots, \bar{x}_j)}\,,
\end{equation}
where the time-ordered product is
\begin{equation}\label{eq:composite-op}
  \hat{M}_j
  = \hat{\Upsilon}(\bar{x}_j)\,U(\tau)\;
    \hat{\Upsilon}(\bar{x}_{j-1})\,U(\tau)
    \;\cdots\;
    \hat{\Upsilon}(\bar{x}_1)\,U(\tau)\,,
\end{equation}
with $U(\tau) = e^{-iH_0\tau}$, and
$p(\bar{x}_1, \ldots, \bar{x}_j)
= \tr[\hat{M}_j\,\rho_S(0)\,\hat{M}_j^\adj]$.
The sequence $(\bar{x}_1, \ldots, \bar{x}_n)$ is the
\textbf{measurement trajectory}.  Averaging over all trajectories
yields the unconditional (non-selective) state
\begin{equation}\label{eq:nonselective-full}
  \rho_S(t_j^+)
  = \int \dd\bar{x}_1 \cdots \dd\bar{x}_j\;
    \hat{M}_j\;\rho_S(0)\;\hat{M}_j^\adj\,.
\end{equation}

% ──────────────────────────────────────────────────────────────
\subsection{The continuous limit}\label{sec:continuous-limit}

Fix a total observation time~$T$ and set $\tau = T/n$.  As
$n \to \infty$, both $\tau \to 0$ and the number of measurements
grows without bound.

\begin{keyresult}[Gaussian integral identity]
  Evaluating~\eqref{eq:nonselective-update} in closed form and
  expanding for large~$\sigma$, one obtains the single-step update
  \begin{equation}\label{eq:single-step-expanded}
    \rho_S(t_j^+)
    = \rho_S(t_j^-)
      - \frac{1}{8\sigma}\bigl[\hat{x}_S,
        [\hat{x}_S, \rho_S(t_j^-)]\bigr]
      + O(1/\sigma^2)\,.
  \end{equation}
\end{keyresult}

\noindent Combining with the free evolution
$\rho_S(t_{j+1}^-) = U(\tau)\,\rho_S(t_j^+)\,U(\tau)^\adj$ and
expanding $U(\tau) = \One - iH_0\tau + O(\tau^2)$, the state
increment over one period is
\begin{equation}\label{eq:state-increment}
  \rho_S(t_{j+1}^+) - \rho_S(t_j^+)
  = -i[H_0, \rho_S]\,\tau
    - \frac{1}{8\sigma}\bigl[\hat{x}_S,
      [\hat{x}_S, \rho_S]\bigr]
    + O(\tau^2)\,.
\end{equation}
Dividing by~$\tau$ and taking $n \to \infty$, the character of the
limiting equation depends on how $\sigma$ scales with~$\tau$:
\begin{enumerate}
  \item \textbf{Quantum Zeno regime} ($\sigma$ fixed):
    the $1/\sigma$ term dominates; the system freezes in a position
    eigenstate.
  \item \textbf{Unitary regime}
    ($\sigma \to \infty$ faster than~$1/\tau$):
    the measurement term vanishes; the system evolves freely.
  \item \textbf{Lindblad regime}
    ($\sigma = 1/(k\tau)$, $k > 0$ fixed):
    both terms contribute at the same order, yielding a non-trivial
    differential equation.
\end{enumerate}

\begin{keyresult}[Master equation from continuous measurement]
  In the Lindblad regime with $\sigma = 1/(k\tau)$ and decoherence
  rate $\Gamma = k/8$, the continuous limit of the non-selective
  evolution is
  \begin{equation}\label{eq:cm-lindblad}
    \eqbox{\frac{\dd\rho_S}{\dd t}
    = -i[H_0,\,\rho_S]
      - \Gamma\,\bigl[\hat{x}_S,
        [\hat{x}_S,\,\rho_S]\bigr]}\,.
  \end{equation}
  This is a Lindblad equation with $L \propto \hat{x}_S$;
  equivalently,
  $-\Gamma[\hat{x}_S,[\hat{x}_S,\rho_S]]
  = 2\Gamma\,\mathcal{D}[\hat{x}_S]\,\rho_S$.
\end{keyresult}

\begin{intuition}[Balancing information and disturbance]
  Each meter in the Lindblad regime is very broad
  ($\sigma \propto 1/\tau \to \infty$) and extracts almost no
  information.  But infinitely many such measurements accumulate a
  finite back-action, producing steady decoherence---precisely the
  mechanism by which a thermal environment induces decoherence.
\end{intuition}

% ──────────────────────────────────────────────────────────────
\section{Quantum stochastic calculus (discrete case)}%
\label{sec:quantum-stochastic-discrete}

\subsection{Free particle under continuous position measurement}%
\label{sec:free-particle}

We illustrate~\eqref{eq:cm-lindblad} for a free particle with
$H_0 = \hat{p}_S^2/(2m)$.

\subsubsection{First moments}

Using the canonical commutation relations
$[\hat{x}_S, \hat{p}_S] = i$, the Ehrenfest equations are
\begin{equation}\label{eq:ehrenfest-free}
  \frac{\dd}{\dd t}\expect{\hat{x}_S}
  = \frac{\expect{\hat{p}_S}}{m}\,,
  \qquad
  \frac{\dd}{\dd t}\expect{\hat{p}_S} = 0\,.
\end{equation}
These are identical to the unitary case: the first moments are
unaffected by the dephasing.

\subsubsection{Momentum diffusion}

The effect of continuous measurement appears in the second moments.
Computing
$\frac{\dd}{\dd t}\expect{\hat{p}_S^2}
= \tr(\hat{p}_S^2\,\dot{\rho}_S)$ yields
\begin{equation}\label{eq:mom-var}
  \eqbox{\frac{\dd}{\dd t}\,\mathrm{Var}(\hat{p}_S) = 2\Gamma}\,.
\end{equation}
The momentum variance grows linearly: continuous position monitoring
induces \textbf{momentum diffusion}.

\begin{keyresult}[Heisenberg uncertainty at work]
  Monitoring position makes it more definite but, by the uncertainty
  principle, makes momentum \emph{less} definite.  The diffusion
  rate~$2\Gamma$ is set by the measurement strength---the same
  parameter that controls the decoherence rate.
\end{keyresult}

\begin{exercise}\label{ex:double-comm-dephasing}
  Verify that $[\hat{x}_S, [\hat{x}_S, \rho_S]]$ in the position
  representation equals $(x - x')^2\,\rho_S(x,x')$, where
  $\rho_S(x,x') = \bra{x}\rho_S\ket{x'}$.  Conclude that diagonal
  elements are unaffected while off-diagonal elements decay at a
  rate proportional to~$(x - x')^2$.
\end{exercise}

\begin{exercise}\label{ex:momentum-diffusion}
  Starting from~\eqref{eq:cm-lindblad} with
  $H_0 = \hat{p}_S^2/(2m)$,
  derive~\eqref{eq:ehrenfest-free} and~\eqref{eq:mom-var}.
  \emph{Hint:}\ use $[\hat{x}_S, \hat{p}_S^n] = in\,\hat{p}_S^{n-1}$
  and the cyclic property of the trace.
\end{exercise}

\begin{exercise}\label{ex:gauss-identity}
  Derive~\eqref{eq:single-step-expanded} by evaluating the
  Gaussian integral in~\eqref{eq:nonselective-update} and expanding
  the result to first order in~$1/\sigma$.
\end{exercise}
